\section{Introduction}
In recent years, generative models have emerged as a cornerstone of machine learning due to their capacity to sample high-dimensional distributions.
Score-based generative models (SGMs) were introduced in their discrete form by Ho et al.~\citep{ho2020denoising} and in their continuous formulation by Song et al.~\citep{song2021scorebased}, building on several years of pioneering work~\citep{sohl2015deep, song2019generative}.
Since then, they have achieved state-of-the-art performance in image and video generation, surpassing GANs, VAEs, and Normalizing Flows \citep{dhariwal2021diffusionbeatGan, muller2022medfusion}. Their frameworks form the foundation of Stable Diffusion \citep{rombach2022high} and many of the most prominent publicly available models today, including DALL·E 3 \citep{Betker2023dalle3}, Imagen 3 \citep{Baldridge2024Imagen3}, and recent advanced architectures like FLUX \citep{blackForestLabs_FLUX1dev_2025}. SGMs have also demonstrated competitive results in audio \citep{chen2020wavegrad, kong2021diffwave} and text generation \citep{li2022diffusionlm, zou2023survey}, positioning them as a flexible and powerful framework across diverse data modalities.

The core idea of SGMs is twofold: the data distribution is gradually corrupted by adding Gaussian noise, and a neural network is trained to reverse this process, recovering the original sample from its corrupted version. New samples are then generated by starting from standard Gaussian noise and iteratively applying the trained network to progressively denoise it.

Theoretical work on SGMs pursues a twofold goal: understanding \textit{how} the approximate backward diffusion converges to the distribution of interest, and proving \textit{that} it does. While it is clear that approximating the score function and using an approximate distribution both introduce errors in generation, it remains an open question under which conditions these approximations are sufficient to guarantee convergence to the true distribution. To establish such guarantees, existing works typically rely on strong assumptions on the data distribution. For instance, \citep{debortoli2022convergence} restricts to bounded support, \citep{gao2025wasserstein} requires log-concavity, and \citep{conforti2023kl} assumes a finite moment of order eight. More broadly, all existing theoretical analyses of SGMs impose some form of strong assumption on the data distribution \citep{bruno2024wasserstein, silveri2025beyond, bortoli2021diffusion, strasman2025analysis}.

A more fundamental issue concerns the backward SDE itself. The founding paper on continuous SGMs \citep{song2021scorebased} relies on \citep{anderson1982reverse} to introduce it, yet \citep{anderson1982reverse} establishes this result only for distributions with bounded support. While other works \citep{bortoli2021diffusion} invoke alternative references \citep{haussmann1986time}, it remains unclear what minimal assumptions on the data distribution are required to rigorously justify the backward SDE, and therefore the diffusion framework itself.

Understanding for which distributions the diffusion framework can be rigorously introduced, and when convergence can be established, is thus a question of both theoretical and practical importance. On the practical side, the recent surge in high-quality image generation has prompted researchers to explore whether diffusion can be extended to heavy-tailed distributions. Several approaches have been proposed, including normalizing flows \citep{hickling2025flexibletails, mcdonald2022cometflows, laszkiewicz2022marginal}, variational autoencoders \citep{lafon2023vae}, GANs \citep{allouche2022evgan, huster2021pareto}, and diffusion-based methods \citep{pandey2025heavytailed, yoon2023levy, shariatian2025heavy}, yet none has fully solved the problem. Several works \citep{pandey2025heavytailed, shariatian2025heavy} explicitly highlight the difficulties SGMs face in this regime, but it remains unclear whether these failures stem from fundamental limitations of the framework or from practical issues such as training instability, poor initialization, or SDE discretization.

In this paper, we show that early stopping and proper initialization together allow the diffusion framework to be applied, theoretically, to any distribution. The idea of improving generation through non-standard initialization has already been explored \citep{zheng2023truncated, zand2024diffusion, DBLP:journals/corr/abs-2306-12360, lee2022priorgrad, shen2025information, scholz2026warm}, but always with the goal of improving efficiency or generation quality. Our contribution is fundamentally different: we provide a theoretical guarantee that diffusion models can be applied to \textit{any} distribution in the sense of KL convergence, without any assumption on the latter.

Precisely, this paper makes the following contributions.
\begin{itemize}
    \item We prove that the backward SDE is well-posed for any data distribution, establishing that the diffusion framework is theoretically grounded in full generality.
    \item We prove that combining early stopping with a flexible initialization scheme that replaces the standard Gaussian prior is sufficient to guarantee convergence of the approximate diffusion to the target distribution without any assumption on the latter, generalizing \cite[Theorem 1]{conforti2023kl} to this end.
    \item We validate our initialization-aware approach on standard benchmarks, demonstrating that a well-chosen initialization improves generation quality while reducing the number of required denoising steps.
    \item We take a first step toward modeling heavy-tailed distributions with standard diffusion models. We discuss what families of distributions are suitable choices for the initialization prior, and provide empirical evidence on toy datasets, synthetic climatic data, and real climatic datasets that proper initialization meaningfully improves over standard diffusion in this regime.
\end{itemize}